\documentclass{situationist-preprint}

\title{The Evidential Inaccessibility of Survival}
\msid{REL-08-b}
\msfamily{REL-08}
\msdiscipline{Philosophy of Religion}
\mstype{Companion manuscript}
\mskeywords{survival of death; personal identity; evidence; replacement; mediumship; confirmation}

\begin{document}
\maketitle

\begin{abstract}
REL-08-a proves that, under a pattern theory of persons, no evidence available before death confirms survival against its matched rivals. We give an argument for the same inaccessibility that does not use the pattern theory. It needs two premises: that what the living can observe depends only on the qualitative states of things that can affect them, and that, on any theory on which survival of death is coherent, a person's surviving is not fixed by the qualitative states of whatever exists after her death. From these it follows that every hypothesis of survival has a replacement rival, on which a qualitatively identical successor exists and the person does not survive, and that no ante-mortem evidence distinguishes the two.
\end{abstract}

\section{The argument}

Let $x$ be a person who dies at $t_d$, and let $S$ be a hypothesis on which $x$ survives: some person existing after $t_d$ is $x$. Evidence is \emph{ante-mortem} for an observer if she acquires it before her own death; the question is whether such evidence could confirm $S$ against hypotheses on which something like $x$ exists after $t_d$ and $x$ does not survive.

\begin{definition}\label{d:rival}
A \emph{replacement rival} of $S$ is a hypothesis $S^{*}$ that agrees with $S$ on the qualitative history of the world up to $t_d$ and on the qualitative states of everything after $t_d$, and on which $x$ does not survive.
\end{definition}

\begin{assumption}[Qualitative locality]\label{a:local}
The probability of any ante-mortem observation, given a hypothesis, depends only on what the hypothesis says about the qualitative states of things that can affect the observer before her death.
\end{assumption}

\begin{assumption}[Non-supervenience]\label{a:nonsup}
If survival of $x$ is coherent, then for every survival hypothesis $S$ there is a coherent replacement rival $S^{*}$.
\end{assumption}

Assumption~\ref{a:nonsup} holds on every theory of persons on which survival of death is a live possibility. On a dualist theory, $S^{*}$ is the hypothesis that $x$'s soul ceases at death and a new soul is created with the same mental states \citep{Swinburne1997}; on a theory of resurrection by reassembly or by divine re-creation, it is the hypothesis that the post-mortem body is a replica rather than $x$ \citep{vanInwagen1978,Hick1976}; on a psychological continuity theory, such as PID-01's, it is the hypothesis that the successor's states, qualitatively the same, were not caused in the way continuity requires. In each case identity across death turns on something, a soul's numerical identity or the right causal route, that is not fixed by the qualitative states of the post-mortem world. A theory on which it were so fixed would make survival consist in the mere existence of a qualitative successor, and would then count every duplicate as the person, which no theory of persons accepts \citep{Parfit1984}.

\begin{theorem}\label{t:inacc}
For every survival hypothesis $S$ and every ante-mortem evidence $E$, there is a coherent rival $S^{*}$ on which $x$ does not survive and $P(E \mid S) = P(E \mid S^{*})$. Hence no ante-mortem evidence confirms $S$ against all rivals on which $x$ does not survive.
\end{theorem}

\begin{proof}
Take $S^{*}$ as in Assumption~\ref{a:nonsup}. By Definition~\ref{d:rival}, $S$ and $S^{*}$ agree on the qualitative states of everything before $t_d$ and after it, and so on the qualitative states of everything that can affect any observer. By Assumption~\ref{a:local}, $P(E \mid S) = P(E \mid S^{*})$. Evidence confirms $S$ against $S^{*}$ only if it is more probable on $S$, so $E$ does not.
\end{proof}

The argument applies equally to evidence from post-mortem communication. If the dead could communicate with the living, Assumption~\ref{a:local} would allow the communications to bear on the hypotheses, but only through their qualitative content, on which $S$ and $S^{*}$ agree. A message carrying information known only to $x$ is evidence that something after $t_d$ has $x$'s memories; it is equally probable whether that thing is $x$ or her replacement \citep{Braude2003}.

\section{Relation to the principal}

Theorem~\ref{t:inacc} is weaker than Theorem~7.7 of REL-08-a in one respect and more general in another. It shows only that some rival is evidentially matched, where REL-08-a identifies two families of matched rivals and derives a ceiling on the probability of survival (Corollary~7.8). It is more general because it does not assume the pattern theory: it holds on any theory meeting Assumption~\ref{a:nonsup}. The family's conclusion therefore does not depend on the details of the pattern theory: survival is coherent under that theory, and the inaccessibility of survival to ante-mortem evidence holds under it and under its rivals.

\begin{objection}
God could reveal to the living that a particular post-mortem person is $x$.
\end{objection}

\begin{reply}
A revelation is an observation with a qualitative content, and by Assumption~\ref{a:local} its probability is the same on $S$ and on $S^{*}$ unless the hypotheses differ about the revealer. If they are specified to differ, the evidence confirms a hypothesis about God's testimony; it bears on survival only through the prior credibility of that testimony, which no ante-mortem evidence can check against the replacement rival.
\end{reply}

\section{Conclusion}

On any theory on which survival of death is coherent, what makes the difference between survival and replacement is not fixed by the qualitative states of things, and what the living observe is. No evidence available before death can confirm survival as such.

\begin{thebibliography}{99}\small
\bibitem[Braude(2003)]{Braude2003} Braude, S.~E. (2003). \emph{Immortal Remains: The Evidence for Life after Death}. Rowman \& Littlefield.
\bibitem[Hick(1976)]{Hick1976} Hick, J. (1976). \emph{Death and Eternal Life}. Collins.
\bibitem[Parfit(1984)]{Parfit1984} Parfit, D. (1984). \emph{Reasons and Persons}. Clarendon Press.
\bibitem[Swinburne(1997)]{Swinburne1997} Swinburne, R. (1997). \emph{The Evolution of the Soul}, revised edn. Clarendon Press.
\bibitem[van Inwagen(1978)]{vanInwagen1978} van Inwagen, P. (1978). The possibility of resurrection. \emph{International Journal for Philosophy of Religion}, 9(2), 114--121.
\end{thebibliography}

\end{document}
